\chapter{State of the Art in Whole-Body Motion Planning for Mobile
  Manipulators}
\label{ch:sota}

% ---------------------------------------------------------------------
\section{Classification of Motion Planning Approaches}
\label{sec:sota-classification}

% TODO: taxonomy -- decoupled vs. whole-body; sampling-based vs.
% optimization-based (trajectory optimization / MPC) vs. reactive
% (potential fields, QP-based inverse kinematics); offline vs. online.
% Cite Sandakalum2022Survey, Thakar2023MobileManip.
\Cref{sec:sota-classification} will classify the main families of motion
planning approaches for mobile manipulators.\footnote{Placeholder --
content to be written.}

% ---------------------------------------------------------------------
\section{Overview of Holistic Motion Planners}
\label{sec:sota-holistic}

% TODO: intro to holistic / whole-body planners.

\subsection{The \ocs{} Framework: Optimal-Control-Based Model Predictive
  Control}
\label{sec:sota-ocs2}

% TODO: OCS2 architecture -- SLQ/DDP solvers, auto-differentiation
% (CppAD), reference manager, MPC-MRT split. Cite ocs2Toolbox,
% Sleiman2021UnifiedMPC.

\subsection{Alternative Frameworks and Methods}
\label{sec:sota-alternatives}

% TODO: e.g. MoveIt/sampling-based, TrajOpt, MPC alternatives
% (acados, control toolbox), reactive QP-IK. Cite moveit2, Mittal2022.

% ---------------------------------------------------------------------
\section{Comparison of Motion Planning Algorithms and Selection of a
  Suitable Algorithm}
\label{sec:sota-selection}

% TODO: comparison table (reactivity, constraint handling, redundancy
% resolution, maturity, ROS 2 support, real-time feasibility) and
% justification for selecting OCS2 for HERo.
